\section{Test-Time Compute Scaling}

\subsection{核心思想}

Training-time scaling（增加训练数据和算力）已被证明有效。Test-Time Compute Scaling 将同样的理念应用于推理阶段：\textbf{在推理时投入更多计算来换取更高质量的输出}。

\subsection{Majority Voting（多数投票）}

最简单的 test-time scaling 方法：

\begin{enumerate}
    \item 对同一个问题（如数学题），让模型生成 $N$ 个不同回答（使用非零温度以获得多样性）
    \item 取出现频率最高的答案作为最终输出
\end{enumerate}

直觉：由于 top-p 采样引入的随机性，模型偶尔会选择错误的 token 导致整条推理链偏离。通过多次采样，正确答案大概率会更频繁地出现。

\subsection{Best-of-N（最优选择）}

Majority Voting 的改进版：

\begin{enumerate}
    \item 生成 $N$ 个不同回答
    \item 使用 \textbf{Reward Model} 或 \textbf{Judge LLM} 对每个回答评分
    \item 选择得分最高的回答
\end{enumerate}

相比 Majority Voting，Best-of-N 利用了外部评估信号，通常效果更好，但需要额外运行一个评估模型。

\subsection{Process Reward Model (PRM)}

更精细的方法——不评估整体答案，而是评估\textbf{每个推理步骤}：

\begin{importantbox}{Process Reward Model 的工作流程}
\begin{enumerate}
    \item LLM 生成部分解答（partial solution）
    \item PRM 对每个步骤评分——分数代表"该步骤最终导向正确答案的概率"
    \item 选择得分最高的步骤继续扩展
    \item 重复上述过程，逐步构建完整解答
\end{enumerate}
这形成了一种\textbf{树搜索}（tree search）机制：在推理步骤的空间中进行搜索，由 PRM 引导搜索方向。
\end{importantbox}

讲者用一个 base-8 乘法的例子演示了这一过程：

\begin{itemize}
    \item Attempt 1：第一步正确，但第二步在十进制而非八进制下计算，PRM 拒绝
    \item Attempt 2：多了一步正确步骤，但最终仍犯同样错误，PRM 再次拒绝
    \item Attempt 4：所有步骤均正确，PRM 通过
\end{itemize}

\subsection{Scaling 实验结果}

Hugging Face 团队的实验显示：

\begin{itemize}
    \item 使用 Llama 3.2 1B 和 3B（小模型），通过增加每题的生成次数
    \item 在数学 benchmark 上，小模型 + test-time scaling 可以\textbf{超越} Llama 3.1 8B 甚至 Llama 3.1 70B
\end{itemize}

\begin{importantbox}{推理速度与输出质量的权衡}
Test-Time Compute Scaling 的本质是用\textbf{推理时间}（inference speed）换\textbf{输出质量}（output quality）。这意味着一个小而高效的模型 + 充足的推理预算，可以在特定任务上匹敌甚至超越体量大得多的模型。这对资源受限的场景具有重要意义。
\end{importantbox}

\begin{knowledgebox}{Reward Model 从何而来}
Test-Time Compute Scaling 中使用的 Reward Model 可以复用 Preference Alignment 阶段（PPO 训练流程中）训练的 reward model。这个模型的输入是文本，输出是 0 到 1 之间的质量分数。也可以直接用 Judge LLM（通过 prompting）替代专门训练的 reward model。
\end{knowledgebox}

\subsection{本章小结}

Test-Time Compute Scaling 是 2025 年初最重要的趋势之一，核心是在推理时投入更多计算以提升输出质量。从简单的 Majority Voting 到 Best-of-N，再到基于 Process Reward Model 的步骤级搜索，方法复杂度和效果依次递增。实验证明小模型 + test-time scaling 可超越大模型。

\newpage
